\documentclass[11pt]{article}
\usepackage[margin=1in]{geometry}
\usepackage{amsmath,amssymb,amsthm,mathtools}
\usepackage[hidelinks]{hyperref}

\newtheorem{theorem}{Theorem}
\newtheorem{lemma}{Lemma}
\newtheorem{remark}{Remark}

\title{An Infinite-Star Homogeneous Space with No Nonconstant Left Quasi-Harmonic Functions}
\author{}
\date{Draft of 18 September 2026}

\begin{document}
\maketitle

\begin{abstract}
Michael Bj\"orklund asked in Appendix II, Problem 2 of \emph{Five remarks about random walks on groups} (arXiv:1406.0763v2) for a countable measured group $(G,\mu)$ and an infinite-index subgroup $H<G$ such that $G/H$ has no nonconstant left (quasi-)$\mu$-harmonic functions. We give an explicit construction using the finitary symmetric group. The induced random walk on the quotient is an infinite star. The leaf equations force every constant-defect solution to grow like the reciprocal of the corresponding transition mass, while integrability at the center forbids a nonzero defect. No boundedness assumption is used. This draft records the source-level argument; publication novelty has not been certified independently.
\end{abstract}

\section{The source question}
Let $G$ be a countable group and let $\mu$ be a probability measure whose support generates $G$ as a semigroup. Following Bj\"orklund, such a pair is a measured group. Left $\mu$-harmonicity is defined using the reflected measure $\check\mu$. In the source, a $\mu$-integrable function $\phi$ is left $\mu$-quasiharmonic if there exists a constant $\ell(\phi)$ such that
\[
\int_G \phi(sg)\,d\check\mu(s)=\phi(g)+\ell(\phi),\qquad g\in G.
\]
Bj\"orklund's Appendix II, \S7.1, Problem 2 asks for an example for which an infinite homogeneous space $G/H$ admits no nonconstant left harmonic or quasiharmonic functions, including unbounded ones.

We prove the required arbitrary-constant-defect statement directly.

\section{Construction}
Let
\[
X=\mathbb N_0=\{0,1,2,\ldots\},\qquad G=\operatorname{Sym}_{\mathrm{fin}}(X),
\]
the group of finitely supported permutations of $X$, and let
\[
H=\operatorname{Stab}_G(0).
\]
For $n\ge1$, write $\tau_n=(0\ n)$. Fix any numbers $p_n>0$ satisfying
\[
\sum_{n=1}^{\infty}p_n=1
\]
and put
\[
\mu(\tau_n)=p_n.
\]
For one completely explicit choice one may take $p_n=2^{-n}$.

\begin{lemma}
The pair $(G,\mu)$ is a countable measured group, $\mu$ is symmetric, and $H$ has infinite index.
\end{lemma}
\begin{proof}
Every finitary permutation has support in $\{0,1,\dots,m\}$ for some $m$, so
\[
G=\bigcup_{m\ge0}\operatorname{Sym}(\{0,1,\dots,m\}),
\]
a countable union of finite groups. Each $\tau_n$ is an involution, hence $\check\mu=\mu$. For distinct $i,j\ge1$,
\[
\tau_i\tau_j\tau_i=(i\ j),
\]
so the support of $\mu$ generates every transposition and therefore every finitary permutation; because the generators are involutions, this is generation as a semigroup as well as as a group.

Finally, the map
\[
G/H\longrightarrow X,\qquad gH\longmapsto g(0)
\]
is a bijection. Indeed, it is well-defined because elements of $H$ fix $0$, and equality of the images implies that $g^{-1}g'$ fixes $0$. Thus $[G:H]=|X|=\infty$.
\end{proof}

\section{The quotient kernel}
Via $G/H\cong X$, left multiplication by $\tau_n$ acts exactly as the permutation $\tau_n$ on $X$. Since $\check\mu=\mu$, the left convolution operator is
\[
(Pu)(x)=\sum_{n\ge1}p_nu(\tau_nx).
\]
Hence
\begin{align}
(Pu)(0)&=\sum_{n\ge1}p_nu(n), \label{eq:center}\\
(Pu)(k)&=p_ku(0)+(1-p_k)u(k),\qquad k\ge1. \label{eq:leaf}
\end{align}
The only infinite tail occurs in \eqref{eq:center}. If $u$ is source-integrable, then this series is absolutely convergent. At a leaf $k$ the corresponding absolute integral is
\[
p_k|u(0)|+(1-p_k)|u(k)|<\infty.
\]

\section{No nonconstant constant-defect functions}
\begin{theorem}\label{thm:main}
Suppose $u:X\to\mathbb C$ is source-integrable at the center,
\[
\sum_{n\ge1}p_n|u(n)|<\infty,
\]
and satisfies
\[
Pu=u+c
\]
for some constant $c\in\mathbb C$. Then $c=0$ and $u$ is constant.
\end{theorem}

\begin{proof}
At a leaf $k\ge1$, equation \eqref{eq:leaf} gives
\[
p_ku(0)+(1-p_k)u(k)=u(k)+c.
\]
Therefore
\[
p_k\bigl(u(0)-u(k)\bigr)=c
\]
and, because $p_k>0$,
\begin{equation}\label{eq:forced}
u(k)=u(0)-\frac{c}{p_k}.
\end{equation}
If $c\ne0$, then
\[
p_k|u(k)|=|p_ku(0)-c|.
\]
Since a positive summable sequence satisfies $p_k\to0$, the right side tends to $|c|>0$. Consequently the summands in $\sum_kp_k|u(k)|$ fail even to tend to zero, contradicting integrability. Hence $c=0$. Equation \eqref{eq:forced} then gives $u(k)=u(0)$ for every $k\ge1$, so $u$ is constant.
\end{proof}

\section{Exact match to Bj\"orklund's quasiharmonic definition}
Let $u:G/H\to\mathbb C$ and let $\widetilde u(g)=u(gH)$ be its right-$H$-invariant pullback. Bj\"orklund's definition says that $\widetilde u$ is left $\mu$-quasiharmonic if it is $\mu$-integrable and there exists a constant $\ell(\widetilde u)$ such that
\[
\int_G\widetilde u(sg)\,d\check\mu(s)
=
\widetilde u(g)+\ell(\widetilde u),\qquad g\in G.
\]
Because $\check\mu=\mu$ and the pullback is right-$H$-invariant, this descends exactly to
\[
Pu=u+\ell(\widetilde u)
\]
on $G/H$. At the basepoint one has
\[
\ell(\widetilde u)=\sum_{n\ge1}p_nu(n)-u(0),
\]
so no normalization $u(0)=0$ is imposed. Theorem~\ref{thm:main}, applied with $c=\ell(\widetilde u)$, gives $\ell(\widetilde u)=0$ and makes $u$ constant. Nonzero constant functions remain allowed, as required: Problem 2 excludes nonconstant functions, not nonzero constants. Ordinary left harmonicity is the special case $c=0$.

\section{Conclusion}
Taking, for instance,
\[
G=\operatorname{Sym}_{\mathrm{fin}}(\mathbb N_0),\qquad
H=\operatorname{Stab}_G(0),\qquad
\mu((0\ n))=2^{-n},
\]
provides an explicit source-valid construction answering Bj\"orklund's Appendix-II Problem 2. In fact every strictly positive probability vector $(p_n)_{n\ge1}$ gives the same conclusion.

\begin{remark}[Prior-art scope]
A bounded literature search dated 18 September 2026 did not locate an explicit later resolution of the exact source problem or this construction. This is negative evidence only and is not a certification of publication novelty. The repository protocol also requires independent audits before this draft can be called audit-complete.
\end{remark}

\begin{thebibliography}{9}
\bibitem{Bjorklund2014}
M.~Bj\"orklund,
\emph{Five remarks about random walks on groups},
arXiv:1406.0763v2, 2014.
\end{thebibliography}

\end{document}
